\begin{abcliste}
\abc
Steht die Kurbel waagrecht, ist der Hebelarm der senkrechten Pedalkraft $F$ die Kurbellänge $r_\mathrm{K}$:
\begin{align}
 M &= \MF \\
   &= \FP\cdot\rK \\
   &= \M \approx \result{\MP{0}{2}}
\end{align}

\abc
Kurbel und Kettenblatt drehen sich zusammen; die Kette zieht am Rand des Kettenblatts und hält dem Drehmoment der Pedalkraft das Gleichgewicht. Ihr Hebelarm ist der Radius $r_\mathrm{B}$ des Kettenblatts:
\begin{align}
 F_\mathrm{Kette} &= \FKF \\
   &= \frac{\M}{\rB} \\
   &= \FK \approx \result{\FKP{0}{2}}
\end{align}
\end{abcliste}
